\documentclass[10pt,a4paper]{amsart}
\usepackage[margin=19mm]{geometry}
\usepackage{amsmath,amssymb,mathtools}
\usepackage[scr=rsfs,cal=euler]{mathalfa}
\usepackage{xfrac}
\usepackage[unicode,hidelinks,bookmarks=false]{hyperref}
\newcommand{\ie}{i.e.\ }
\newcommand{\cf}{cf.\ }
\newcommand{\resp}{respectively\ }
\newcommand{\Subsection}{\subsection}
\providecommand{\nobreakdash}{\nobreak\hbox{-}}
\setlength{\emergencystretch}{2em}
\multlinegap0pt
\usepackage{mathtools}
\usepackage{amssymb}
\usepackage[shortlabels]{enumitem}
\newtheorem{proposition}{Proposition}
\title{A note on explicit homological invariants of graded double Ore extensions}
\author{Andr\'es Rubiano}
\date{Source: arXiv:2604.02456v1 (CC BY 4.0)}
\begin{document}
\maketitle

\noindent\emph{Source and licence.} A note on explicit homological invariants of graded double Ore extensions. Version arXiv:2604.02456v1 (CC BY 4.0), distributed by arXiv under the Creative Commons Attribution 4.0 International licence, http://creativecommons.org/licenses/by/4.0/, as recorded in the arXiv licence metadata for this exact version and shown on the abstract page https://arxiv.org/abs/2604.02456v1. Full licence notice: \texttt{licenses/arxiv-2604.02456-CC-BY-4.0.txt}.\par\smallskip

\noindent\textit{Editorial scope.} Counted unit: the article's proposition prop:Mx (source0.tex lines 535--553). The article's complete original proof of it is delivered with it (source0.tex lines 555--575, one \texttt{proof} environment of 21 lines). No mathematics is written, repaired or re-derived: the statement and the proof are the article's own lines, in the article's own notation, and no retained line is substituted or re-flowed.  Retained uncounted as the source-local context the proof needs: lines 288--295.  Nothing was omitted from any retained span, so the package carries no omission record.  No retained line contains a citation, so the wrapper carries no bibliography and no bibliography entry.  No retained line contains a figure environment, an image-inclusion command or drawing code, so no image is delivered and the package declares no asset dependency.  Editorial formatting only, in addition: an AMS \texttt{amsart} preamble that restates the article's own declarations and macros used by the retained lines, namely the environments \texttt{proposition} on separate counters, together with the standard packages those lines need.  The retained environments print in this document as Proposition 1 in order of appearance, whereas the article prints the same items with the numbering of its own full document.  The complete native source of the article as distributed in the arXiv e-print for this exact version is bundled unchanged as \texttt{sources/197736/source0.tex}.
\begin{align}
x_2x_1 &= -x_1x_2,\label{eq:K1}\\
y_2y_1 &= -y_1y_2,\label{eq:K2}\\
y_1x_1 &= x_1y_1,\label{eq:K3}\\
y_1x_2 &= x_2y_2,\label{eq:K4}\\
y_2x_1 &= x_1y_2,\label{eq:K5}\\
y_2x_2 &= \alpha x_2y_1.\label{eq:K6}
\end{align}

\begin{proposition}\label{prop:Mx}
The sequence
\[
0\longrightarrow B(-2)\xrightarrow{\,d_2^{(x)}\,}B(-1)^2\xrightarrow{\,d_1^{(x)}\,}B\longrightarrow M_x\longrightarrow 0,
\]
where
\[
d_1^{(x)}(u,v)=ux_1+vx_2,
\
d_2^{(x)}(w)=(wx_2,wx_1),
\]
is a minimal graded free resolution of \(M_x\). In particular,
\[
\beta_{0,0}^B(M_x)=1,\
\beta_{1,1}^B(M_x)=2,\
\beta_{2,2}^B(M_x)=1,
\]
and all other graded Betti numbers vanish.
\end{proposition}

\begin{proof}
Clearly \(\operatorname{coker} d_1^{(x)}\cong M_x\) and
\[
d_1^{(x)}d_2^{(x)}(w)=w(x_2x_1+x_1x_2)=0
\]
by \eqref{eq:K1}, so \(\operatorname{Im}d_2^{(x)}\subseteq \operatorname{Ker} d_1^{(x)}\). The map \(d_2^{(x)}\) is injective because \(B\) is a domain. From the exact sequence
\[
0\longrightarrow \operatorname{Ker} d_1^{(x)}\longrightarrow B(-1)^2\xrightarrow{\,d_1^{(x)}\,}B\longrightarrow M_x\longrightarrow 0
\]
and the formulas \(H_B(t)=1/(1-t)^4\), \(H_{M_x}(t)=1/(1-t)^2\), we obtain
\[
H_{\operatorname{Ker} d_1^{(x)}}(t)
=
2tH_B(t)-H_B(t)+H_{M_x}(t)
=
\frac{t^2}{(1-t)^4}
=
H_{B(-2)}(t).
\]
Therefore \(\operatorname{Ker} d_1^{(x)}=\operatorname{Im}d_2^{(x)}\) and the sequence is exact. Minimality is immediate.
\end{proof}

\end{document}
